\documentclass[11pt]{article}
\usepackage[T1]{fontenc}
\usepackage{lmodern,microtype,amsmath,amssymb,amsthm,booktabs}
\microtypesetup{expansion=false}
\usepackage[margin=1.05in,headheight=14pt]{geometry}
\usepackage{xcolor}
\definecolor{linknavy}{RGB}{25,55,125}
\usepackage[colorlinks=true,linkcolor=linknavy,citecolor=linknavy,urlcolor=linknavy]{hyperref}
\usepackage{orcidlink,fancyhdr}
\hypersetup{
  pdftitle={A stability refinement for simple critical zeros in short intervals},
  pdfauthor={Felipe Santibanez-Leal},
  pdfsubject={Short-interval zeta-zero proportions and finite Hilbert-space stability},
  pdfkeywords={Riemann zeta function, simple zeros, short intervals, pair correlation, Gram matrix, interval arithmetic}
}
\newtheorem{theorem}{Theorem}[section]
\newtheorem{proposition}[theorem]{Proposition}
\newtheorem{lemma}[theorem]{Lemma}
\newtheorem{corollary}[theorem]{Corollary}
\theoremstyle{remark}
\newtheorem{remark}[theorem]{Remark}
\DeclareMathOperator{\tr}{tr}
\DeclareMathOperator{\supp}{supp}
\newcommand{\R}{\mathbb R}
\newcommand{\C}{\mathbb C}
\newcommand{\HS}{\mathrm{HS}}
\newcommand{\cE}{\mathcal E}
\newcommand{\cC}{\mathcal C}
\newcommand{\versiondoi}{10.5281/zenodo.22727389}
\newcommand{\conceptdoi}{10.5281/zenodo.22727388}
\newcommand{\recordurl}{https://github.com/fsantibanezleal/CAOS_RESEARCH/tree/main/problems/number-theory/riemann-hypothesis}
\pagestyle{fancy}
\fancyhf{}
\fancyhead[L]{\small CAOS Research Preprint}
\fancyhead[R]{\small Short-interval stability\textperiodcentered\ v0.01}
\fancyfoot[C]{\thepage}
\fancypagestyle{plain}{\fancyhf{}\fancyfoot[C]{\thepage}\renewcommand{\headrulewidth}{0pt}}
\allowdisplaybreaks[1]

\title{A stability refinement for simple critical zeros\\in short intervals}
\author{Felipe Santiba\~nez-Leal\,\orcidlink{0000-0002-0150-3246}\\
\small CAOS Research program\\
\small\texttt{github.com/fsantibanezleal/CAOS\_RESEARCH}}
\date{2026-09-12\\[2pt]\normalsize Version 0.01}

\begin{document}
\maketitle
\thispagestyle{plain}
\begin{center}\small
\fbox{\parbox{0.92\textwidth}{\centering
\textbf{PREPRINT}\textperiodcentered\ not peer reviewed\textperiodcentered\ CAOS Research\\
Riemann-hypothesis research series, first paper\\
Version 0.01\textperiodcentered\ 2026-09-12\textperiodcentered\ CC BY 4.0\\
Version DOI: \href{https://doi.org/\versiondoi}{\versiondoi}\\
Concept DOI: \href{https://doi.org/\conceptdoi}{\conceptdoi} (latest version)\\
\href{\recordurl}{Public sources, proof audit, and reproducible certificates}
}}
\end{center}

\begin{abstract}
For each fixed $\theta$ for which Wang's short-interval cosine bound
$c(\theta)=2-\theta/2-\cot(\theta/\sqrt2)/\sqrt2$ is positive, we give an
explicit constant strictly larger than $c(\theta)$ as a lower asymptotic
proportion of simple critical zeros in $(T,T+T^\theta]$. The corresponding
distinct-zero bound also improves. We transfer the known stable rank--trace
inequality in the \texttt{ainta} manuscript directly to Lamzouri's finite
Hilbert operator, retain the Gram defect of its simple real points, and
combine three shifted partitions into consecutive triples. A quantitative
obstruction to additive triples of cosine-kernel zeros supplies a positive
gap for every fixed admissible $\theta$. Wang's fixed-test pair-correlation
theorem provides the arithmetic input; smoothing is performed after the
large-height limit. An interval certificate at $\theta=3/4$ improves the
simple-critical constant from $0.4190750129754243\ldots$ to
$0.4190768284253039\ldots$, and gives distinct proportion
$0.7095384142126519\ldots$. The general strict improvement has an analytic
proof independent of numerical minimization. This does not improve the
positivity exponent, claim a global record, or resolve the Riemann hypothesis.
\end{abstract}

\section{Statement and relation to prior work}

Let $N(T,H)$ count the nontrivial zeros $\rho=\beta+i\gamma$ of $\zeta(s)$
with $T<\gamma\le T+H$, including multiplicity. Let $N_0^s(T,H)$ count the
simple zeros among them with $\beta=1/2$, and let $N^d(T,H)$ count distinct
complex zeros, each once. Every proportion below has denominator $N(T,H)$.

The recent argument discovered by Claude and communicated by Alp\"oge and
Furman~\cite{AF} made the optimized Montgomery--Taylor simple-critical
proportion $0.672500703679\ldots$ unconditional. Lamzouri~\cite{Lamzouri}
replaced its finite-compression machinery by a Hilbert-space inequality for
conjugation-invariant multisets. Wang~\cite{Wang} established the needed
short-interval pair-correlation theorem and deduced, for fixed $0<\theta<1$,
\begin{equation}\label{eq:wang}
 \liminf_{T\to\infty}\frac{N_0^s(T,T^\theta)}{N(T,T^\theta)}\ge c(\theta),
 \qquad
 \liminf_{T\to\infty}\frac{N^d(T,T^\theta)}{N(T,T^\theta)}
 \ge\frac{1+c(\theta)}2,
\end{equation}
where
\begin{equation}\label{eq:c}
 c(\theta)=2-\frac\theta2-\frac1{\sqrt2}\cot(\theta/\sqrt2).
\end{equation}
The derivative is $c'(\theta)=\tfrac12\cot^2(\theta/\sqrt2)>0$. Since
$c(\theta)\to-\infty$ as $\theta\downarrow0$ and $c(1)>0$, there is a
unique zero $\theta_0\in(0,1)$, numerically
$\theta_0=0.550193964744154\ldots$.

The stable rank--trace inequality, convex Gram defect, pinching method, and
cosine-root non-additivity used here are prior work in the
\texttt{ainta} repository~\cite{Ainta}. Further global window and block
refinements are also under active development~\cite{Trmdy}. We do not claim
these ingredients or an improved global record. Our contribution is their
direct finite-Hilbert transfer, an explicit quantitative three-point bound,
and the strict refinement of Wang's entire positive short-interval curve.
The finite proof below is included to make the treatment of off-line zeros,
multiplicities, and singular Gram matrices explicit.

\begin{theorem}[Strict short-interval refinement, D]\label{thm:main}
Fix $\theta_0<\theta<1$ and $R>2/c(\theta)$. Define
\begin{equation}\label{eq:explicit}
\begin{aligned}
 a&=\theta/\sqrt2,\qquad q=a\tan a,\qquad M=\pi\theta R,\\
 B&=q(1+M^2/a^2),\\
 b&=\min\left\{\frac q2,
       \frac{q^4}{4(q+M)^2(M+3q/2)}\right\},\qquad
 \delta=\frac12\min\{1,2(b/B)^2\}.
\end{aligned}
\end{equation}
Then $0<\delta\le1/2$, and, with
\begin{equation}\label{eq:cstar}
 c_*(\theta,R)=c(\theta)
       +\frac{\delta\{c(\theta)-2/R\}}{3-\delta}>c(\theta),
\end{equation}
we have
\begin{equation}\label{eq:mainbounds}
 \liminf_{T\to\infty}\frac{N_0^s(T,T^\theta)}{N(T,T^\theta)}
 \ge c_*(\theta,R),\qquad
 \liminf_{T\to\infty}\frac{N^d(T,T^\theta)}{N(T,T^\theta)}
 \ge\frac{1+c_*(\theta,R)}2.
\end{equation}
In particular, $R=4/c(\theta)$ gives a fully specified positive improvement.
\end{theorem}

The labels [D] and [MV] distinguish mathematical deductions from finite
machine-verifiable statements. The conservative gap in
Theorem~\ref{thm:main} requires no numerical search. A stronger finite
certificate gives the following example.

\begin{corollary}[Certified example, D+MV]\label{cor:example}
Set $\theta=3/4$, $R=21/4$, $d=1/7000$, and
\[
 C_*:=\frac{3c(3/4)-2d/R}{3-d}
 =0.419076828425303996736665787527\ldots.
\]
Then the simple-critical proportion in \eqref{eq:mainbounds} is at least
$C_*$, and the distinct-zero proportion is at least
\[
 \frac{1+C_*}2=0.709538414212651998368332893763\ldots.
\]
The improvement in the simple-critical proportion is
$0.000001815449879663002112176828\ldots$.
\end{corollary}

The conclusion is for each fixed exponent and all sufficiently large
heights in the asymptotic sense. We do not provide an effective starting
height or uniformity as $\theta\downarrow\theta_0$. In particular, this
does not lower Wang's positive-proportion threshold. The analytic dependency
is Wang's inspected preprint theorem, whose underlying pair-correlation
method is due to Baluyot, Goldston, Suriajaya, and
Turnage-Butterbaugh~\cite{BGSTB,BGSTBcorrection}.

\section{A stable inequality for finite multisets}

Let $Z$ be a finite conjugation-invariant multiset in $\C$. We sum over its
distinct support with positive integer multiplicities $m_z=m_{\bar z}$,
and put $N=\sum_zm_z$. Choose a real even compactly supported
$\eta\in L^2(\R)$ with $\int\eta^2=1$, and define
\[
 v_z(u)=\eta(u)e^{-2\pi izu},\qquad
 K(t)=\widehat{\eta^2}(t),\qquad
 \widehat f(t)=\int_{\R}f(u)e^{-2\pi itu}\,du.
\]
Every $v_z$ belongs to $L^2$ because $Z$ is finite and the support is compact.
Work in the finite-dimensional span $E$ of these vectors. Our inner product
is $\langle v,w\rangle=\int\bar v w$, linear in the second variable; write
$|v\rangle\langle w|$ for $h\mapsto v\langle w,h\rangle$.

Define the reflected rank-one operator
\begin{equation}\label{eq:operator}
 A=\sum_zm_z|v_z\rangle\langle v_{\bar z}|.
\end{equation}
It is selfadjoint by conjugation invariance. Direct multiplication gives
\begin{equation}\label{eq:trace}
 \tr A=N,\qquad
 Q:=\|A\|_{\HS}^2=\tr A^2
   =\sum_{z,w}m_zm_wK(z-w)^2.
\end{equation}
Indeed, $\langle v_{\bar z},v_z\rangle=\int\eta^2=1$, while the trace
of the product of the $z$ and $w$ terms, before multiplicity factors, is
$K(w-z)K(z-w)=K(z-w)^2$, since $K$ is even. The ordinary square in
\eqref{eq:trace} is essential. Individual terms for complex arguments need
not be positive or even real; their complete sum is a nonnegative operator
norm.

Let $x_1<\cdots<x_s$ be the simple real points of $Z$. Let $r$ be the
number of distinct multiple real points and $k$ the number of distinct
nonreal conjugate pairs. Put
\[
 P=\sum_{j=1}^s|v_{x_j}\rangle\langle v_{x_j}|=VV^*,\qquad
 G=V^*V=(K(x_i-x_j))_{i,j=1}^s,\qquad C=A-P.
\]
Then $G$ is positive semidefinite with diagonal one and $\tr P=s$.
For each nonreal pair put
$g=(v_z+v_{\bar z})/2$ and $h=(v_z-v_{\bar z})/(2i)$. Its contribution
to $C$ is $2m_z(|g\rangle\langle g|-|h\rangle\langle h|)$.
The form of $C$ is nonpositive on the orthogonal complement of the $r$
multiple-real vectors and $k$ vectors $g$. Thus
\begin{equation}\label{eq:inertia}
 n_+(C)\le r+k,\qquad N\ge s+2r+2k.
\end{equation}
No distance of a nonreal point from the real axis is assumed.

\begin{lemma}[Known stable rank--trace inequality]\label{lem:stable}
Suppose $P=VV^*$ on a finite-dimensional Hilbert space, the $s$ columns of
$V$ have norm one, $C$ is selfadjoint, and $n_+(C)\le d$. Define
\[
 \Psi(t)=\begin{cases}(t-1)^2,&0\le t\le2,\\2t-3,&t\ge2.
 \end{cases}
\]
Then
\begin{equation}\label{eq:stable}
 \|P+C\|_{\HS}^2\ge4\tr(P+C)-3s-4d+\tr\Psi(V^*V).
\end{equation}
\end{lemma}
\begin{proof}
This is the unit-column case of the stability inequality in
\cite[Section 2]{Ainta}. We give a spectral proof that keeps all zero
eigenvalues. Let $m=\dim E$ and let $p_1\le\cdots\le p_m$ be the
eigenvalues of $P$. If $d<m$, the min-max principle gives
$\lambda_i(P+C)\le p_{i+d}$ for $1\le i\le m-d$: intersect the span
of the first $i+d$ eigenvectors of $P$ with a codimension-at-most-$d$
subspace on which $C\le0$.

It follows that
\[
 \tr(P+C-2I)^2\ge\sum_{j=d+1}^m(2-p_j)_+^2
 \ge\sum_{j=1}^m(2-p_j)_+^2-4d.
\]
Here $0\le(2-p_j)_+^2\le4$. If $d\ge m$, the final right-hand side is
nonpositive and the same inequality is immediate. For every $p\ge0$,
$(2-p)_+^2=\Psi(p)-2p+3$. The nonzero spectra of $VV^*$ and $V^*V$
agree, and $\Psi(0)=1$, so
\[
 \tr\Psi(P)=\tr\Psi(V^*V)+m-s.
\]
Substitute these identities and $\tr P=s$, and expand the left-hand
square. All dimensional terms cancel and give \eqref{eq:stable}.
\end{proof}

\begin{proposition}[Finite Hilbert stability transfer]\label{prop:finite}
For the multiset and operator above, write $D(G)=\tr\Psi(G)$ and
$D_Z=s+r+2k$ for the number of distinct support points. Then
\begin{equation}\label{eq:finite}
 s\ge2N-Q+D(G),\qquad
 D_Z\ge\frac{3N-Q+D(G)}2.
\end{equation}
\end{proposition}
\begin{proof}
By Lemma~\ref{lem:stable} and \eqref{eq:inertia},
$Q\ge4N-3s-4(r+k)+D(G)$. The multiplicity bound in
\eqref{eq:inertia} makes this at least $2N-s+D(G)$. For the second
claim, the difference between $4N-3s-4(r+k)$ and $3N-2D_Z$ is
$N-s-2r\ge0$.
\end{proof}

This proof retains the signed off-line remainder throughout. The positive
Gram defect belongs only to the actual simple real points. The second
inequality in \eqref{eq:finite}, rather than elementary counting from a
simple-zero proportion, will yield the distinct-zero companion.

\section{From three-point energy to a density gain}

For a positive semidefinite matrix $M$ with diagonal one, define
$\cE(M)=2\sum_{i<j}|M_{ij}|^2$.

\begin{lemma}[Block defect and pinching]\label{lem:pinch}
We have $D(M)\ge\min\{1,\cE(M)\}$. If $M_j$ are the principal blocks
of a partition of $M$, then $D(M)\ge\sum_jD(M_j)$.
\end{lemma}
\begin{proof}
If all eigenvalues of $M$ are at most 2, then
$D(M)=\tr(M-I)^2=\cE(M)$. If some eigenvalue exceeds 2, its own
$\Psi$ value exceeds 1 and all other terms are nonnegative.
For pinching, diagonalize each principal block and let $(w_i)$ be the
resulting orthonormal basis of the full space. Scalar Jensen applied to
the spectral resolution of $M$ gives
\[
 \sum_jD(M_j)=\sum_i\Psi(\langle w_i,Mw_i\rangle)
 \le\sum_i\langle w_i,\Psi(M)w_i\rangle=D(M).
\]
Only scalar convexity of $\Psi$ is used; no operator-convex assertion is
needed.
\end{proof}

\begin{lemma}[Three shifted partitions]\label{lem:triples}
Suppose $0<d\le1$, $R>0$, and the real even kernel satisfies
\begin{equation}\label{eq:triangle}
 2\{K(u)^2+K(v)^2+K(u+v)^2\}\ge d,
 \qquad u,v\ge0,\quad u+v\le R.
\end{equation}
If $\ell=x_s-x_1$ is the span of the simple real points, with $\ell=0$
when $s\le1$, then
\begin{equation}\label{eq:packing}
 D(G)\ge\frac d3\left(s-2-\frac{2\ell}{R}\right).
\end{equation}
\end{lemma}
\begin{proof}
For $s\ge3$, the sum of the spans of its $s-2$ consecutive triples is
at most $2\ell$, since each successive gap occurs at most twice. Hence
at most $2\ell/R$ triples have span greater than $R$. Every remaining
triple has defect at least $d$ by Lemma~\ref{lem:pinch}.

Separate the starting indices of consecutive triples into their three
residue classes modulo 3. Within each class the triples are disjoint
principal blocks; add the unused points as singleton blocks. Apply
pinching to each partition and average the three inequalities. Every
short triple has been counted once, proving \eqref{eq:packing}. If its
right-hand side is negative, nonnegativity proves the inequality; this
also covers $s<3$.
\end{proof}

The factor $1/3$ records the three shifted partitions. Overlapping triples
cannot simply be summed without this step.

\section{A quantitative obstruction for the cosine kernel}

For $0<\lambda\le1$ put
\begin{equation}\label{eq:density}
 f_\lambda(u)=\frac{\cos(\sqrt2u)}{\sqrt2\sin(\lambda/\sqrt2)}
       \mathbf1_{[-\lambda/2,\lambda/2]}(u),\qquad
 K_\lambda=\widehat f_\lambda.
\end{equation}
This is a nonnegative probability density. Set $a=\lambda/\sqrt2$,
$q=a\tan a$, and $X=\pi\lambda t$. Integration yields
\begin{equation}\label{eq:kernel}
 K_\lambda(t)=
 \frac{\cos X-(X/a)\cot(a)\sin X}{1-X^2/a^2}.
\end{equation}
The apparent singularities at $X=\pm a$ are removable. Equivalently,
with $\operatorname{sinc}(x)=\sin(x)/x$ continuously extended at zero,
\begin{equation}\label{eq:sinc}
 K_\lambda(t)=
 \frac{\operatorname{sinc}(X-a)+\operatorname{sinc}(X+a)}
      {2\operatorname{sinc}(a)}.
\end{equation}

An actual positive root satisfies $X\tan X=q$. If $X,Y,X+Y$ all
corresponded to positive roots, tangent addition would force
$X^2+XY+Y^2+q^2=0$. The following estimate gives a quantitative version
without dividing by a possibly vanishing cosine or removable denominator.

\begin{lemma}[Explicit positive triangle energy]\label{lem:gap}
Fix $R>0$, put $M=\pi\lambda R$, and define $B,b$ as in
\eqref{eq:explicit}, using $a=\lambda/\sqrt2$ and $q=a\tan a$.
Then, throughout $u,v\ge0$, $u+v\le R$,
\begin{equation}\label{eq:gap}
 2\{K_\lambda(u)^2+K_\lambda(v)^2+K_\lambda(u+v)^2\}
 \ge2(b/B)^2>0.
\end{equation}
\end{lemma}
\begin{proof}
Let $A_X=q\cos X-X\sin X$. For $Z=X+Y$, angle addition gives
\begin{equation}\label{eq:addition}
 (X^2+XY+Y^2+q^2)\sin X\sin Y
 =A_XA_Y-XA_X\sin Y-YA_Y\sin X-qA_Z.
\end{equation}
Take $X,Y\ge0$, $Z\le M$, and put
$h=\max\{|A_X|,|A_Y|,|A_Z|\}$. If $h\le q/2$, then
\[
 q\le q|\cos X|+q|\sin X|
 \le h+(q+X)|\sin X|,
\]
so $|\sin X|\ge q/[2(q+M)]$, and likewise for $Y$. Taking absolute
values in \eqref{eq:addition} therefore yields
\[
 \frac{q^4}{4(q+M)^2}\le h^2+(M+q)h
 \le(M+3q/2)h.
\]
If $h>q/2$, the other alternative in the definition of $b$ applies.
In either case $h\ge b>0$. The identity
\[
 A_X=q(1-X^2/a^2)K_\lambda(X/(\pi\lambda))
\]
holds also at the removable points by continuity. Since
$|q(1-X^2/a^2)|\le B$ for $0\le X\le M$, at least one of the three
kernel values has absolute value at least $b/B$. This proves
\eqref{eq:gap}.
\end{proof}

The zero-set obstruction itself is already used in~\cite{Ainta}. The
explicit formula in Lemma~\ref{lem:gap} makes its positive compact gap
available for every fixed exponent without a numerical minimum. Its
small size is acceptable for the qualitative strict-improvement theorem;
Section~\ref{sec:certificate} certifies a more useful gap in one example.

\section{Short-interval arithmetic and passage to the sharp density}

We record the exact analytic interface to avoid a support or smoothing
uniformity assumption. Fix $0<\lambda<\theta<1$, put $H=T^\theta$,
$L=\log T$, and $I=(T,T+H]$. Wang's Theorem 2.2 states that for every
fixed real even $g\in C_c^\infty(\R)$ supported in $[-\lambda,\lambda]$,
\begin{align}\label{eq:wangtest}
 W_I(g)&:=\sum_{\gamma,\gamma'\in I}
 \widehat g\!\left(\frac{i(\rho-\rho')L}{2\pi}\right)
 \frac4{4-(\rho-\rho')^2}\notag\\
 &=\frac{HL}{2\pi}\left(g(0)+\int_\R|\alpha|g(\alpha)d\alpha\right)
      +O_g(H+T^\lambda L^2).
\end{align}
The constants may depend on fixed $\lambda,\theta$. This is an
unconditional complex-zero formula; RH is not an assumption.

Choose a fixed real even $\eta\in C_c^\infty((-\lambda/2,\lambda/2))$
with $\int\eta^2=1$, and write $f=\eta^2$, $K=\widehat f$, $g=f*f$.
Both $g$ and $g''$ meet the fixed-test hypotheses. At
$z=i(\rho-\rho')L/(2\pi)$, differentiation of the Fourier transform gives
\[
 \widehat {g''}(z)=-4\pi^2z^2K(z)^2
                  =L^2(\rho-\rho')^2K(z)^2.
\]
The exact identity
\[
 \left(\widehat g(z)-\frac{\widehat {g''}(z)}{4L^2}\right)
 \frac4{4-(\rho-\rho')^2}=K(z)^2
\]
removes the rational weight. Apply \eqref{eq:wangtest} separately to
the two fixed tests and combine the resulting estimates. As in
\cite[Lemma 3.1]{Wang}, this gives
\begin{equation}\label{eq:unweighted}
 \sum_{\gamma,\gamma'\in I}K(z_\rho-z_{\rho'})^2
 =\cC(f)\frac{HL}{2\pi}+O_f(H+T^\lambda L^2),
\end{equation}
where
\begin{equation}\label{eq:functional}
 z_\rho=\frac{i(\rho-1/2)L}{2\pi},\qquad
 \cC(f)=\int f^2+\iint|u-v|f(u)f(v)\,du\,dv.
\end{equation}
Indeed, the main term for $g''$ is
$-\int(f')^2+2\int f^2$, finite for the fixed smooth density, and its
coefficient is $1/(4L^2)$. A varying scalar coefficient after applying
the theorem to two fixed tests does not require a varying-test theorem.

Reflection $\rho\mapsto1-\bar\rho$ preserves the height interval and
multiplicity and maps $z_\rho$ to $\bar z_\rho$. Thus the multiset
$Z_T=\{z_\rho:\gamma\in I\}$ satisfies Proposition~\ref{prop:finite}.
Its simple real points are precisely the simple critical zeros. Their
real-coordinate span is at most $X_T=HL/(2\pi)$. The short-interval
Riemann--von Mangoldt formula gives
\begin{equation}\label{eq:normalization}
 N(T,H)=X_T+O(H+\log T),\qquad X_T/N(T,H)\longrightarrow1.
\end{equation}
After division by $N$, the error in \eqref{eq:unweighted} is
$O_f(1/L+T^{\lambda-\theta}L)=o(1)$. The strict support gap
$\lambda<\theta$ is therefore necessary at this stage.

The cosine density \eqref{eq:density} is Wang's minimizing density and
satisfies
\begin{equation}\label{eq:cost}
 \cC(f_\lambda)=\frac\lambda2+
         \frac1{\sqrt2}\cot(\lambda/\sqrt2)=2-c(\lambda).
\end{equation}
It is not a smooth compactly supported test because its endpoint values
do not vanish. Choose even cutoffs $0\le\chi_\varepsilon\le1$, supported
strictly inside $(-\lambda/2,\lambda/2)$ and approaching one in the
interior. Normalizing $\eta_\varepsilon=\chi_\varepsilon\sqrt{f_\lambda}/
\|\chi_\varepsilon\sqrt{f_\lambda}\|_2$ gives densities
$f_\varepsilon=\eta_\varepsilon^2$ converging in both $L^1$ and $L^2$
to $f_\lambda$. Moreover $f_\lambda\to f_\theta$ in those spaces as
$\lambda\uparrow\theta$, and $\cC$ is continuous for these convergences
on bounded support~\cite[Section 4]{Wang}.

For normalized nonnegative densities,
\begin{equation}\label{eq:uniform}
 \sup_{x\in\R}|\widehat f(x)-K_\theta(x)|\le\|f-f_\theta\|_1,
 \qquad |\widehat f(x)|,|K_\theta(x)|\le1.
\end{equation}
Thus the doubled triple energy changes by at most
$12\|f-f_\theta\|_1$. This retains a fixed positive energy margin
through the smooth approximation.

\begin{proposition}[Transfer with a certified limiting gap]\label{prop:transfer}
Fix $0<\theta<1$, $R>0$, and $0<d\le1$. Suppose \eqref{eq:triangle}
holds for $K_\theta$. Define
\[
 c_d=\frac{c(\theta)-2d/(3R)}{1-d/3}.
\]
Then the simple-critical and distinct proportions in \eqref{eq:mainbounds}
are at least $c_d$ and $(1+c_d)/2$, respectively.
\end{proposition}
\begin{proof}
First fix $0<d'<d$. Choose a sequence of fixed smooth densities $f_j$
supported in $(-\lambda_j/2,\lambda_j/2)$, with
$\lambda_j<\theta$, converging to $f_\theta$ in $L^1$ and $L^2$.
By \eqref{eq:uniform}, every sufficiently large $j$ has triple energy
at least $d'$. Fix such a $j$ before letting $T$ grow.

Write $s=N_0^s(T,H)$, $N=N(T,H)$, and $Q$ for the pair sum with this
fixed kernel. Propositions~\ref{prop:finite} and Lemma~\ref{lem:triples},
with actual span at most $X_T$, imply
\[
 \left(1-\frac{d'}3\right)\frac sN
 \ge2-\frac QN-\frac{2d'}{3N}-\frac{2d'X_T}{3RN}.
\]
Let $T\to\infty$, using \eqref{eq:unweighted} and
\eqref{eq:normalization}, and then let $j\to\infty$. The result is
$\liminf s/N\ge c_{d'}$.

For the distinct count, the other finite inequality gives
\[
 \liminf\frac{N^d(T,H)}N
 \ge\frac{1+c(\theta)}2+\frac{d'}6(c_{d'}-2/R)
 =\frac{1+c_{d'}}2.
\]
The equality uses $c_{d'}-c(\theta)=d'(c_{d'}-2/R)/3$.
Finally let $d'\uparrow d$; all scalar expressions are continuous since
$d\le1<3$. No test function changes inside a large-$T$ limit.
\end{proof}

\begin{proof}[Proof of Theorem~\ref{thm:main}]
Lemma~\ref{lem:gap} with $\lambda=\theta$ proves a limiting-kernel gap
at least $2(b/B)^2$, and hence at least the positive $\delta$ in
\eqref{eq:explicit}. Apply Proposition~\ref{prop:transfer} with
$d=\delta$. Its constant equals \eqref{eq:cstar}. The gain is strictly
positive because $R>2/c(\theta)$.
\end{proof}

\section{The certified numerical example}\label{sec:certificate}

For $\theta=3/4$ and $R=21/4$, the recorded finite computation proves
\begin{equation}\label{eq:numericalgap}
 2\{K_{3/4}(u)^2+K_{3/4}(v)^2+K_{3/4}(u+v)^2\}\ge\frac1{7000}
 \quad(u,v\ge0,\ u+v\le21/4).
\end{equation}
The certificate is a complete binary partition of the square $[0,R]^2$.
Every endpoint and threshold is rational. A box is either subdivided,
proved wholly outside the triangle, or accepted using an energy lower
bound valid on the entire box. The final tree contains 48,761 nodes,
24,252 accepted energy leaves, 129 outside leaves, and no unresolved
boxes~\cite{Record}.

For completeness, the analytic enclosure used at each leaf is simple.
Since $f_\theta$ is a probability density on $[-\theta/2,\theta/2]$,
\begin{equation}\label{eq:lipschitz}
 |K_\theta'(x)|\le2\pi\int|u|f_\theta(u)du\le\pi\theta.
\end{equation}
On a real interval with center $x_0$ and radius $r$, an outward interval
evaluation of $K_\theta(x_0)$ therefore supplies the valid lower bound
\[
 |K_\theta(x)|\ge
 \max\{0,\operatorname{lower}(|K_\theta(x_0)|)-\pi\theta r\}.
\]
Apply this separately to the box ranges for $u$, $v$, and $u+v$, square
the nonnegative lower bounds, and double their sum. Every acceptance
comparison requires the entire lower enclosure to exceed the rational
threshold. An outside box is discarded only if the sum of its two lower
coordinates is strictly greater than $R$, so the triangle boundary is
retained.

Discovery used Arb through \texttt{python-flint==0.9.0} at 160-bit
precision and the removable-singularity-free formula \eqref{eq:sinc}.
Each accepted leaf was replayed at 256 bits with a separately derived
Taylor evaluator. For real $x$,
\begin{equation}\label{eq:taylor}
 \operatorname{sinc}(x)=\sum_{j=0}^{95}\frac{(-1)^jx^{2j}}{(2j+1)!}
 +\mathcal R(x),\qquad
 |\mathcal R(x)|\le\frac{|x|^{192}}{193!}.
\end{equation}
The remainder follows by applying the real sine Taylor theorem through
degree 192 and dividing by $x$, with continuous extension at zero.
Arb encloses both the polynomial and its remainder. This replay is an
independent kernel evaluator; it shares the partition reconstruction,
Lipschitz argument, and Arb library with the first computation. It is
not a second independent full verifier or a Lean proof.

\begin{proof}[Proof of Corollary~\ref{cor:example}]
Use the finite certificate \eqref{eq:numericalgap} in
Proposition~\ref{prop:transfer}, with $d=1/7000$ and $R=21/4$.
The full threshold $d$ is available because the proposition first uses
every $d'<d$ through smoothing and then passes to $d$. Substitution
gives the stated constants. The outward rational enclosures for their
values and positive difference are recorded with the certificate.
\end{proof}

\begin{table}[ht]
\centering
\begin{tabular}{@{}lll@{}}
\toprule
Proportion at $\theta=3/4$ & Wang baseline & This certificate\\
\midrule
Simple and critical & $0.4190750129754243\ldots$ & $0.4190768284253039\ldots$\\
Distinct & $0.7095375064877121\ldots$ & $0.7095384142126519\ldots$\\
\bottomrule
\end{tabular}
\caption{All denominators count nontrivial zeros with multiplicity in
$(T,T+T^{3/4}]$. The entries are proportions, not percentages.}
\end{table}

The certificate's role is quantitative. Even without it,
Theorem~\ref{thm:main} proves a strict improvement for every fixed
$\theta>\theta_0$ using its explicit analytic gap. Conversely, the
finite numerical inequality alone does not prove a theorem about zeta
zeros without the finite operator and fixed-test analytic transfer.

\appendix
\section{Reproducibility, provenance, and scope}

The public record~\cite{Record}, experiment
\texttt{EXP-002-short-interval-stability}, contains the mathematical proof,
an adversarial audit, the interval runner, the canonical partition tree,
and rational enclosures for the constants. The hypothesis and permitted
parameter family were declared before computation in repository commit
\texttt{266486f}. A
deterministic floating grid selected candidate radii; its sampled values
were never used as certified lower bounds. The final certification used
CPU arithmetic; no GPU was necessary.

From the repository root, with the problem's pinned requirements installed,
a fresh output directory preserves the canonical artifacts:
\begin{verbatim}
python problems/number-theory/riemann-hypothesis/experiments/\
EXP-002-short-interval-stability/run.py \
  --output-dir tmp/riemann-replay --radius 21/4 --threshold 1/7000
pytest tests/test_riemann_certificates.py
\end{verbatim}
The displayed backslashes indicate line continuation for readability;
use one command line on shells with different continuation syntax.
The runner fails with an exact resumable checkpoint if its time or node
budget expires while boxes remain. It does not report a certificate in
that state. A separate bounded trial exercised termination and exact
resume behavior.

The tree replay checks its hash, reconstructs every box, verifies all
accepted leaves, and rejects missing, trailing, unknown, or miscounted
nodes. A hash protects accidental corruption; geometric replay is still
required when a modified tree has a recomputed hash. Adversarial checks
rejected false whole-square acceptance, a false outside-domain claim,
truncation, appended data, altered counts, an unresolved flag, an invalid
exponent, and an inflated threshold. Nineteen focused tests, including
one complete 256-bit certificate replay, passed in the recorded run.
These are implementation checks, not formal verification of the entire
analytic proof.

The canonical preorder tree has SHA-256 digest
\begin{center}\small\ttfamily
78829635026aeacd69c02213cc899b5029ef2566bbfaf1c86b757cc5f92ffe55
\end{center}

The source comparison was completed on 2026-09-12. It used the full
Wang v1 manuscript, Lamzouri v2, the current Alp\"oge--Furman paper,
and pinned snapshots of the directly relevant global stability
repositories. Queries combining the Wang identifier, short intervals,
stability, Gram defects, and exact numerical signatures found no prior
public statement of this short-interval refinement. This supports the
scoped contribution claimed here, without establishing exhaustive
priority or detecting unpublished concurrent work. The known finite
stability inequality and cosine-root mechanism are explicitly credited
to the repository identifier \texttt{ainta}; no personal author identity
is inferred from that identifier.

The theorem relies on established classical analytic inputs through
Wang's recent preprint, not on an end-to-end machine formalization of
those inputs. The analytic and numerical reviews can still contain
errors, and the two numerical evaluators share arithmetic and geometric
components. No effective height threshold, uniform exponent limit,
lower positivity threshold, global-record claim, or elimination of all
off-line zeros follows from this work.

AI assistance was used for derivation, source research, adversarial
review, and implementation. These automated reviews do not constitute
independent peer review. The new manuscript and original research
exposition are released under CC BY 4.0; code follows the repository's
MIT license. Cited third-party materials
retain their own licenses. In particular, local copies of arXiv texts
with nonexclusive-distribution licenses are not relicensed by this
manuscript.

\begin{thebibliography}{9}\small\raggedright

\bibitem{AF}
L.~Alp\"oge and R.~Furman,
\emph{More than two thirds of the zeta zeros are simple and on the
critical line}, arXiv:2608.13637v2 (2026).
Proof discovered by Claude and verified and communicated by the listed
authors, as described in the paper's provenance.
\url{https://arxiv.org/abs/2608.13637v2}.

\bibitem{Lamzouri}
Y.~Lamzouri,
\emph{A new proof that more than $2/3$ of the zeros of the Riemann zeta
function are simple and on the critical line}, arXiv:2609.02882v2 (2026).
\url{https://arxiv.org/abs/2609.02882v2}.

\bibitem{Wang}
B.~Wang,
\emph{Simple critical zeros and distinct zeros of the Riemann
zeta-function in short intervals}, arXiv:2609.07918v1 (2026).
\url{https://arxiv.org/abs/2609.07918v1}.

\bibitem{Ainta}
\texttt{ainta} (repository maintainer identifier),
\emph{More than 67.3\% of the zeros of the Riemann zeta function are
simple and lie on the critical line}, manuscript in
\texttt{zeta-simple-zeros}, 2026, MIT license.
Pinned commit \texttt{040c5e899e658aed7b56a2a87f501798fe10761d},
\texttt{paper/riemann.tex}; accessed 2026-09-12.
\href{https://github.com/ainta/zeta-simple-zeros/tree/040c5e899e658aed7b56a2a87f501798fe10761d}{Pinned manuscript and source snapshot on GitHub}.

\bibitem{Trmdy}
\texttt{trmdy} repository and contributors,
\emph{Zeta simple-zero proportion refinement}, research repository
\texttt{zeta-simple-zeros-673137}, 2026, MIT license.
Pinned commit \texttt{1610b97b7895ff34982260f8dcaf04a0f7b82cf7};
accessed 2026-09-12. Candidate global bounds are not used as analytic
premises of this paper.
\href{https://github.com/trmdy/zeta-simple-zeros-673137/tree/1610b97b7895ff34982260f8dcaf04a0f7b82cf7}{Pinned contributor manuscript and certificates on GitHub}.

\bibitem{BGSTB}
S.~A.~C.~Baluyot, D.~A.~Goldston, A.~I.~Suriajaya, and C.~L.~Turnage-Butterbaugh,
\emph{An unconditional Montgomery theorem for pair correlation of zeros
of the Riemann zeta-function}, Acta Arith. \textbf{214} (2024), 357--376;
arXiv:2306.04799.
\url{https://arxiv.org/abs/2306.04799}.

\bibitem{BGSTBcorrection}
S.~A.~C.~Baluyot, D.~A.~Goldston, A.~I.~Suriajaya, and C.~L.~Turnage-Butterbaugh,
\emph{Pair correlation of zeros of the Riemann zeta function I:
Proportions of simple zeros and critical zeros},
arXiv:2501.14545v3 (2026), especially Section 3 on the
corrected error term in the earlier Montgomery theorem.
\url{https://arxiv.org/abs/2501.14545v3}.

\bibitem{Record}
F.~Santiba\~nez-Leal,
\emph{Riemann-hypothesis research record, EXP-002: short-interval
stability}, CAOS Research, 2026. Proof, source comparisons, code, exact
partition certificate, and adversarial audit.
\url{https://github.com/fsantibanezleal/CAOS_RESEARCH/tree/main/problems/number-theory/riemann-hypothesis}.

\end{thebibliography}
\end{document}
